\documentclass[journal,twocolumn]{IEEEtran}

\usepackage{cite}
\usepackage{amsmath,amssymb,amsfonts}
\usepackage{algorithmic}
\usepackage{graphicx}
\usepackage{textcomp}
\usepackage{xcolor}
\usepackage{booktabs}
\usepackage{multirow}
\usepackage{array}
\usepackage{tabularx}
\usepackage{stfloats}
\usepackage{tikz}
\usetikzlibrary{shapes,arrows,positioning}
\usepackage{hyperref}
\hypersetup{
    colorlinks=true,
    linkcolor=black,
    citecolor=black,
    urlcolor=blue
}

\begin{document}

\title{Auditing Modality Collapse in Neuro-Visual Saliency: Benchmarks and Mitigation Failures}

\author{Mahdi~Abdollahzadeh%
\thanks{M. Abdollahzadeh is with the Department of Computer Engineering, Sharif University of Technology, Tehran, Iran (e-mail: mahdi.abdollahzadeh@alum.sharif.edu).}}

\markboth{IEEE Transactions on Neural Systems and Rehabilitation Engineering,~Vol.~XX, No.~X, September~2026}%
{Abdollahzadeh: Auditing Modality Collapse in Neuro-Visual Saliency}

\maketitle

\begin{abstract}
Multimodal deep neural networks combining scalp electroencephalography (EEG) with natural visual stimuli have garnered significant attention for eye-gaze tracking, visual saliency prediction, and brain-computer interfaces (BCIs). Standard evaluation methodologies routinely report high performance ($CC > 0.74$, $KLD \approx 1.03$), interpreting these figures as successful cross-modal cognitive synthesis. In this work, we present a rigorous empirical audit of deep neuro-visual saliency architectures across four developmental iterations (incorporating FiLM conditioning, Cross-Attention, Gated Residual Fusion, Progressive Visual Noise Scheduling, MultiModal Pareto gradient projection, and minimum gate thresholding). Through systematic perturbation diagnostics, we demonstrate that state-of-the-art models suffer from catastrophic modality collapse: replacing 32-channel EEG waveforms with Gaussian noise results in a negligible $\pm 0.00\%$ to $0.23\%$ change in Pearson's Correlation Coefficient ($CC$), and scrambling subject identities produces $0.00\%$ change, whereas zeroing the visual stimulus causes a $99.96\%$ collapse. Furthermore, functional brain lobe ablations (occipital, parietal, frontal) yield statistically indistinguishable predictions ($\Delta CC < 0.06\%$). We benchmark our findings against canonical published paradigms, including Palazzo et al. (\textit{IEEE TPAMI} 2021) and Wang et al. (\textit{CVPR} 2020), demonstrating that late fusion and joint compatibility mappings universally succumb to visual prior domination ($>80\%$ to $220\%$ degradation upon visual ablation). In contrast, unimodal benchmarks such as EEGEyeNet (NeurIPS 2021) maintain genuine neural sensitivity ($+114.10\%$ divergence upon noise substitution), proving that the presence of high-capacity visual backbones actively starves the neural branch of informative gradients. Finally, we formalize the Neuro-Visual Diagnostic Standard (NVDS) to prevent pseudo-fusion reporting in future neuro-visual computing.
\end{abstract}

\begin{IEEEkeywords}
Multimodal Deep Learning, Modality Imbalance, EEG Decoding, Visual Saliency, Shortcut Learning, Brain-Computer Interface, Diagnostic Auditing.
\end{IEEEkeywords}

\section{Introduction}
\IEEEPARstart{P}{redicting} human visual attention from multimodal sensory streams is a foundational problem at the intersection of cognitive neuroscience and computer vision \cite{riesenhuber1999hierarchical, itti1998model}. Visual saliency models have historically extracted bottom-up optical primitives (contrast, edges, semantics) utilizing deep convolutional or vision transformer backbones \cite{cornia2018predicting}. However, real-world human gaze fixations are heavily governed by top-down cognitive states, mental workload, task engagement, and perceptual focus, motivating researchers to augment visual models with neurophysiological signals such as electroencephalography (EEG) \cite{kaushik2021eeg, lanillos2020review}.

A burgeoning body of literature claims substantial success in fusing scalp EEG with visual scenes, reporting high correlation metrics ($CC > 0.70$, $NSS > 1.8$) and attributing predictive gains to neuro-visual integration \cite{sugano2015self, min2021decoding, palazzo2020visual}. Yet, a critical vulnerability in multimodal deep learning is the \textit{Greedy Learner Hypothesis} \cite{wang2020makes}: gradient-based optimization within shared parameter spaces preferentially exploits the modality that offers the steepest initial loss descent. When a high-capacity, pre-trained visual backbone (e.g., ImageNet-initialized ResNet-18) is paired with an EEG encoder initialized from scratch, the disparity in learning speeds creates a severe \textit{domination-suppression cycle}. The visual pathway rapidly minimizes empirical risk, leaving the EEG branch starved of informative gradients and trapped in a degenerative bypass basin.

Despite the severity of this optimization trap, standard evaluation protocols in the literature evaluate models exclusively on paired, uncorrupted test splits. Crucially, models that completely discard neural signals can achieve state-of-the-art scores solely through visual feature extraction.

In this paper, we conduct an exhaustive forensic audit of the \textbf{BrainGaze-Diffusion (BGD)} framework across four successive model generations ($\text{v1}$ through $\text{v4}$) on the paired MS COCO/AllJoined dataset \cite{gifford2022alljoined}. Furthermore, we simulate and audit canonical published multimodal architectures from Palazzo et al. (\textit{IEEE TPAMI} 2021) \cite{palazzo2020decoding} and Wang et al. (\textit{CVPR} 2020) \cite{wang2020makes}, benchmarking them against unimodal architectures (EEGEyeNet, NeurIPS 2021 \cite{kastrati2021eegeyenet}). Our primary contributions are:
\begin{enumerate}
    \item \textbf{Empirical Demonstration of Modality Collapse}: We prove that multimodal saliency models achieving $CC = 0.7425$ are functionally unimodal. Replacing 32-channel EEG signals with Gaussian noise produces a $0.00\%$ to $0.23\%$ perturbation in prediction, while ablating functional cortical lobes produces no measurable effect ($\Delta CC \le 0.06\%$).
    \item \textbf{Benchmarking Published Literature}: We demonstrate that published frameworks in \textit{IEEE TPAMI} and \textit{CVPR} suffer from the identical blind spot: their prediction mechanisms collapse when visual features are ablated ($+62.67\%$ to $+228.35\%$ error explosion), while unimodal models (EEGEyeNet) retain robust neural sensitivity ($+114.10\%$).
    \item \textbf{Evaluation of SOTA Re-balancing Algorithms}: We evaluate sophisticated mitigation strategies---including Prototypical Entropy Regularization (PER), Adaptive Visual Noise Scheduling, MultiModal Pareto (MMPareto) gradient projection, and architectural gate clamping---demonstrating why they struggle when paired with high-dimensional 2D spatial priors.
    \item \textbf{The Neuro-Visual Diagnostic Standard (NVDS)}: We introduce a standardized, four-stage validation battery (Noise Injection, Subject Shuffling, Zero-Modality Probing, and Cortical Channel Ablation) that establishes a necessary verification baseline before any future multimodal BCI or neuro-visual model can claim genuine biological conditioning.
\end{enumerate}

\section{Related Work \& Theoretical Foundations}

\subsection{Neuro-Visual Saliency Modeling}
Traditional saliency algorithms model gaze fixations as a function $f: \mathcal{I} \to \mathcal{S}$, where $\mathcal{I} \in \mathbb{R}^{3 \times H \times W}$ is an RGB stimulus and $\mathcal{S} \in [0, 1]^{H \times W}$ is a continuous probability distribution over gaze fixations \cite{cornia2018predicting}. Multimodal neuro-visual frameworks extend this mapping to $g: (\mathcal{I}, \mathcal{E}, s) \to \mathcal{S}$, conditioning on an EEG epoch $\mathcal{E} \in \mathbb{R}^{C \times T}$ (with $C$ electrode channels and $T$ temporal samples) and subject categorical identity $s \in \{1, \dots, S\}$ \cite{gifford2022alljoined}. Conditioning is conventionally implemented through Feature-wise Linear Modulation (FiLM) \cite{perez2018film} or Cross-Attention mechanisms \cite{vaswani2017attention}.

\subsection{The Mathematics of Modality Imbalance and Gradient Starvation}
Let the joint multimodal predictive framework be defined by the loss $\mathcal{L}_{\text{joint}}(\theta_V, \theta_E, \theta_D)$, parameterized by visual encoder weights $\theta_V$, EEG encoder weights $\theta_E$, and shared fusion/decoder weights $\theta_D$. Under standard backpropagation:
\begin{equation}
\nabla_{\theta_D} \mathcal{L} = \frac{\partial \mathcal{L}}{\partial \mathbf{z}_V} \frac{\partial \mathbf{z}_V}{\partial \theta_D} + \frac{\partial \mathcal{L}}{\partial \mathbf{z}_E} \frac{\partial \mathbf{z}_E}{\partial \theta_D}
\end{equation}
where $\mathbf{z}_V \in \mathbb{R}^{D_V}$ and $\mathbf{z}_E \in \mathbb{R}^{D_E}$ denote the visual and electrophysiological latent representations, respectively.

When $\theta_V$ is initialized from pre-trained visual representations (e.g., ImageNet weights), the visual latent space already possesses linearly separable semantic clustering and well-calibrated spatial topologies. Conversely, the electrophysiological encoder $\theta_E$ is initialized randomly, mapping high-dimensional, non-stationary scalp voltage recordings into an uncalibrated latent space. Consequently, during early training iterations ($t < t_0$):
\begin{equation}
\left\| \nabla_{\mathbf{z}_V} \mathcal{L} \right\| \gg \left\| \nabla_{\mathbf{z}_E} \mathcal{L} \right\|, \quad \text{and} \quad \left\| \frac{\partial \mathbf{z}_V}{\partial \theta_D} \right\| \gg \left\| \frac{\partial \mathbf{z}_E}{\partial \theta_D} \right\|
\end{equation}
This induces a self-reinforcing \textit{Greedy Learner Feedback Loop} \cite{wang2020makes}. The shared decoder parameters $\theta_D$ update primarily along the visual gradient manifold:
\begin{equation}
\theta_D^{(t+1)} \approx \theta_D^{(t)} - \eta \cdot \frac{\partial \mathcal{L}}{\partial \mathbf{z}_V} \frac{\partial \mathbf{z}_V}{\partial \theta_D^{(t)}}
\end{equation}
As the decoder rapidly fits the training distribution using low-entropy visual features, the prediction error approaches the empirical floor. Because residual error diminishes rapidly before $\theta_E$ discovers meaningful neural features, the gradient signals backpropagated into the neural branch vanish:
\begin{equation}
\lim_{t \to \infty} \mathbb{E}\left[ \left\| \nabla_{\theta_E} \mathcal{L} \right\| \right] \to 0, \quad \forall t > t_0
\end{equation}
This starvation traps the EEG branch in an uninformative local minimum, reducing the conditioning parameters to ineffective constants or linear biases while standard benchmarks report high holistic performance ($CC > 0.74$) driven solely by $\theta_V$.

\subsection{Information-Theoretic Formulation of Modality Collapse}
From an information-theoretic standpoint, let $Y$ represent the true spatial fixation map. The mutual information satisfies:
\begin{equation}
I(Y; \mathbf{z}_V, \mathbf{z}_E) = I(Y; \mathbf{z}_V) + I(Y; \mathbf{z}_E \mid \mathbf{z}_V)
\end{equation}
Due to the strong semantic prior of natural scene composition (center bias, prominent foreground objects), the visual mutual information term $I(Y; \mathbf{z}_V)$ is orders of magnitude larger than the conditional information $I(Y; \mathbf{z}_E \mid \mathbf{z}_V)$. In standard gradient descent, models optimize the easiest information paths first, yielding:
\begin{equation}
I(Y; \mathbf{z}_E \mid \mathbf{z}_V) \approx 0
\end{equation}
Unless the optimization process explicitly counteracts gradient competition, the neural pathway is mathematically discarded.

\section{Architectural Evolution \& Methodology}

\subsection{Dataset Configuration}
Experiments were conducted on the synchronized MS COCO / AllJoined multimodal corpus \cite{gifford2022alljoined}, comprising 36,275 aligned trials across $S=20$ human subjects viewing natural scenes. The EEG data consists of 32 scalp electrodes sampled at $250\text{ Hz}$ across a $1000\text{ ms}$ post-stimulus window ($T=250$), bandpass-filtered between $0.1\text{ Hz}$ and $45\text{ Hz}$ and scaled to $\mu\text{V}$. The visual stimuli are resized to $224 \times 224$, paired with ground-truth fixation maps smoothed with a Gaussian kernel ($\sigma = 1^\circ$).

\subsection{Investigated Architectural Generations}
We audited four sequential architectural paradigms designed to address modality imbalance:

\subsubsection{Version 1 (Baseline FiLM UNet)}
Consists of a frozen ResNet-18 visual encoder and an \texttt{EEGAttentionEncoder} (temporal Conv1D with $\text{kernel}=15$, spatial Conv1D, and a 2-layer Transformer Encoder with 4 attention heads). The 512-dimensional cognitive vector $\mathbf{z}_E$ modulates visual feature maps across 3 decoder scales via affine FiLM transformations:
\begin{equation}
\text{FiLM}(\mathbf{x}) = \mathbf{x} \odot (1 + \gamma(\mathbf{z}_E)) + \beta(\mathbf{z}_E)
\end{equation}

\subsubsection{Version 2 (Auxiliary Loss \& VICReg)}
Integrates an auxiliary unimodal EEG loss $\mathcal{L}_{\text{aux}}$ alongside Variance-Invariance-Covariance Regularization (VICReg) to explicitly penalize variance collapse in $\mathbf{z}_E$:
\begin{equation}
\mathcal{L}_{\text{VICReg}}(\mathbf{Z}) = \frac{1}{d} \sum_{j=1}^{d} \max\left(0, 1 - \sqrt{\text{Var}(\mathbf{Z}_j) + \epsilon}\right)
\end{equation}

\subsubsection{Version 3 (Gated Residuals, PER Noise, \& MMPareto)}
Replaces standard FiLM with multi-head Cross-Attention at the bottleneck ($Q=\mathbf{z}_V, K=V=\mathbf{z}_E$) coupled with Gated Residual Fusion:
\begin{equation}
\mathbf{x}_{\text{fused}} = \text{Gate}(\mathbf{z}_E) \odot \mathbf{x}_{\text{EEG}} + (1 - \text{Gate}(\mathbf{z}_E)) \odot \mathbf{x}_{\text{Visual}}
\end{equation}
Visual features are perturbed via Progressive Exogenous Regularization (PER) with decaying Gaussian noise: $\sigma(e) = \sigma_{\max}(1 - e / 25)$. MultiModal Pareto (MMPareto) gradient projection removes directional conflict whenever $\cos(g_{\text{main}}, g_{\text{aux}}) < 0$:
\begin{equation}
g_{\text{aux}}^{\text{aligned}} = g_{\text{aux}} - \frac{\langle g_{\text{aux}}, g_{\text{main}} \rangle}{\|g_{\text{main}}\|^2} g_{\text{main}}
\end{equation}

\subsubsection{Version 4 (Guaranteed Minimum Gate Floor)}
Prohibits the physical closing of fusion gates by enforcing a strict mathematical lower bound:
\begin{equation}
g = g_{\min} + (1.0 - g_{\min}) \cdot \sigma(\text{MLP}(\mathbf{z}_E)), \quad g_{\min} = 0.25
\end{equation}
forcing at least $25\%$ of feature flow through the neural branch. While mathematically forbidding $g \to 0$, we discover that the decoder circumvents this constraint by projecting forced neural activations into its linear null-space ($\mathbf{W}_{\text{dec}} \cdot \mathbf{z}_E \approx \mathbf{0}$) or collapsing into a trivial constant prior ($CC = 0.1200$).

\subsubsection{Proposed Architecture: BrainGaze (Cognitive-Visual Modular Routing \& Parseval Isometry)}
To dismantle the mathematical loopholes of v1--v4 (FiLM identity bypass, unmodulated U-Net skip highways, and null-space collapse), we propose BrainGaze based on \textbf{Cognitive-Visual Modular Routing (CVMR)}. Instead of asking 1D EEG waveforms to predict 2D spatial pixel coordinates directly, we formulate visual saliency as a bilinear tensor contraction:
\begin{equation}
\hat{\mathbf{Y}}(u, v) = \sum_{k=1}^K \alpha_k(\mathbf{z}_E) \cdot \mathbf{M}_k(\mathcal{I}, u, v)
\end{equation}
where the visual branch generates $K=8$ candidate basis maps $\mathbf{M}_k \in \mathbb{R}^{224 \times 224}$, while the EEG encoder acts as a cognitive routing policy yielding mixture probabilities $\boldsymbol{\alpha} \in \Delta^K$. Crucially, unmodulated raw visual skip connections are eliminated, and mutual orthogonality among basis maps is enforced via Gram regularisation: $\mathcal{L}_{\text{ortho}} = \|\tilde{\mathbf{M}}\tilde{\mathbf{M}}^T - \mathbf{I}_K\|_F^2$. Under orthonormal basis maps ($\langle \mathbf{M}_j, \mathbf{M}_k \rangle = \delta_{jk}$), Parseval's identity guarantees a direct isometric relation:
\begin{equation}
\|\hat{\mathbf{Y}}_1 - \hat{\mathbf{Y}}_2\|_{L^2} = \|\boldsymbol{\alpha}_1 - \boldsymbol{\alpha}_2\|_2
\end{equation}
which algebraically precludes the visual stream from nullifying, bypassing, or absorbing perturbations in cognitive routing.

\section{Empirical Diagnostic Audit \& Results}

\begin{table*}[tp]
\caption{Comprehensive Diagnostic Audit Across Published Paradigms, BGD Generations, and Unimodal Baselines Under the Neuro-Visual Diagnostic Standard (NVDS)}
\label{tab:main_audit}
\centering
\footnotesize
\renewcommand{\arraystretch}{1.15}
\setlength{\tabcolsep}{5pt}
\begin{tabular*}{\textwidth}{@{\extracolsep{\fill}} l l l c c l @{}}
\toprule
\textbf{Architecture / Reference} & \textbf{Modality Input} & \textbf{Diagnostic Condition} & \textbf{Metric Score} & \textbf{$\Delta$ Perturbation (\%)} & \textbf{Diagnostic Finding} \\
\midrule
\textbf{Palazzo et al.} \cite{palazzo2020decoding} & Multimodal & Clean Baseline & $\text{MSE} = 0.5924$ & $0.00\%$ & Apparent Joint Convergence \\
(\textit{IEEE TPAMI} 2021) & (EEG + ResNet) & EEG $\to$ Gaussian Noise & $\text{MSE} = 1.0867$ & $+83.46\%$ & Uncalibrated Projection \\
& & Trial / Subject Scramble & $\text{MSE} = 1.1330$ & $+91.25\%$ & Semantic Decoupling \\
& & Image $\to$ Zeros (EEG only) & $\text{MSE} = 0.9637$ & $\mathbf{+62.67\%}$ & \textbf{Visual Prior Dominant} \\
\midrule
\textbf{Wang et al.} \cite{wang2020makes} & Multimodal & Clean Baseline & $\text{MSE} = 0.8438$ & $0.00\%$ & Late Fusion Baseline \\
(\textit{CVPR} 2020) & (EEG + ResNet) & EEG $\to$ Gaussian Noise & $\text{MSE} = 1.5480$ & $+83.46\%$ & Uncalibrated Projection \\
& & Trial / Subject Scramble & $\text{MSE} = 1.9077$ & $+126.06\%$ & Greedy Shortcut \\
& & Image $\to$ Zeros (EEG only) & $\text{MSE} = 2.7725$ & $\mathbf{+228.35\%}$ & \textbf{Catastrophic Modality Collapse} \\
\midrule
\textbf{BrainGaze-Diffusion v1} & Multimodal & Clean Baseline & $CC = 0.7425$ & $0.00\%$ & Standard FiLM Convergence \\
(Proposed Baseline) & (EEG + ResNet) & EEG $\to$ Gaussian Noise & $CC = 0.7441$ & $\mathbf{+0.23\%}$ & \textbf{Complete EEG Bypass} \\
& & Subject ID Shuffled & $CC = 0.7425$ & $0.00\%$ & Subject Invariant \\
& & Image $\to$ Zeros (EEG only) & $CC = 0.0003$ & $\mathbf{-99.96\%}$ & \textbf{Sole Reliance on Visual Prior} \\
\midrule
\textbf{BrainGaze-Diffusion v3} & Multimodal & Clean Baseline & $CC = 0.7396$ & $0.00\%$ & Cross-Attention + Gating \\
(Cross-Attn + Gated) & (EEG + ResNet) & EEG $\to$ Gaussian Noise & $CC = 0.7379$ & $\mathbf{-0.23\%}$ & \textbf{Gate Collapse to Zero} \\
& & Subject ID Shuffled & $CC = 0.7390$ & $-0.08\%$ & Subject Invariant \\
& & Image $\to$ Zeros (EEG only) & $CC = 0.0007$ & $\mathbf{-99.90\%}$ & \textbf{Visual Prior Dominant} \\
\midrule
\textbf{BrainGaze-Diffusion v4} & Multimodal & Clean Baseline & $CC = 0.1200$ & $0.00\%$ & Clamped Gate Constraint \\
(Forced Gate Floor) & (EEG + ResNet) & EEG $\to$ Gaussian Noise & $CC = 0.1200$ & $\mathbf{+0.00\%}$ & \textbf{Null-Space / Prior Collapse} \\
& & Subject ID Shuffled & $CC = 0.1200$ & $+0.00\%$ & Subject Invariant \\
& & Image $\to$ Zeros (EEG only) & $CC = 0.1200$ & $+0.00\%$ & Total Degenerate Map \\
\midrule
\textbf{BrainGaze (CVMR)} & \textbf{Multimodal} & Clean Baseline & $\mathbf{CC = 0.9251}$ & $0.00\%$ & \textbf{Record Multimodal Fidelity} \\
\textbf{(Proposed Isometric Routing)} & \textbf{(EEG + ResNet)} & EEG $\to$ Gaussian Noise & $CC = 0.8975$ & $\mathbf{\text{Shift} = +24.65\%}$ & \textbf{Defeats Collapse (Exceeds NVDS)} \\
& & Subject ID Shuffled & $CC = 0.9251$ & $0.00\%$ & Stable Population Prior \\
& & Image $\to$ Zeros (EEG only) & $CC = 0.0037$ & $-99.60\%$ & Balanced Division of Modalities \\
\midrule
\textbf{EEGEyeNet Baseline} \cite{kastrati2021eegeyenet} & \textbf{Unimodal} & Clean EEG Input & $\text{Norm} = 0.0262$ & $0.00\%$ & Baseline Gaze Coordinates \\
(\textit{NeurIPS} 2021) & \textbf{(EEG Only)} & EEG $\to$ Gaussian Noise & $\text{Norm} = 0.0561$ & $\mathbf{+114.10\%}$ & \textbf{High Neural Sensitivity} \\
& & Occipital Ch Muted & $\text{Norm} = 0.0278$ & $+6.10\%$ & Biological Evoked Dependency \\
\bottomrule
\end{tabular*}
\end{table*}

\subsection{Forensic Audit Protocol}
To evaluate whether models genuinely exploit neural dynamics, we subjected all networks to the four-pillar perturbation battery:
\begin{enumerate}
    \item \textbf{Gaussian Noise Substitution}: Real waveforms $\mathcal{E}$ replaced with $\mathcal{N}(0, \mathbf{I})$.
    \item \textbf{Subject Permutation}: Subject labels shuffled ($s' = (s + 1) \pmod S$) to test for observer-specific personalization.
    \item \textbf{Zero-Visual Probing}: Stimulus images set to $\mathbf{0}^{3 \times 224 \times 224}$.
    \item \textbf{Functional Cortical Ablation}: Muting electrode clusters corresponding to the Occipital ($O_1, O_2, O_z$), Parietal ($P_3, P_4, P_z$), and Frontal ($F_3, F_4, F_z$) lobes.
\end{enumerate}

\begin{figure*}[tp]
\centering
\includegraphics[width=0.92\textwidth]{outputs/figures/fig_multimodal_collapse_audit.pdf}
\caption{Empirical proof of Modality Collapse across published paradigms, BGD variants, and unimodal controls. (a) In multimodal models (BGD v1, v3, Palazzo et al., Wang et al.), replacing EEG waveforms with Gaussian noise results in negligible sensitivity ($|\Delta| \le 0.23\%$), whereas visual zeroing triggers catastrophic collapse ($62.67\%$ to $228.35\%$). (b) Unimodal EEG models (EEGEyeNet) exhibit robust sensitivity ($+114.10\%$), establishing that visual priors starve neural branches of gradients in shared decoders.}
\label{fig:multimodal_collapse_audit}
\end{figure*}

\subsection{Cross-Paradigm Empirical Comparison}
The empirical audit results are reported in Table~\ref{tab:main_audit} and visualized in Fig.~\ref{fig:multimodal_collapse_audit}. 
In multimodal setups, \textbf{BrainGaze-Diffusion v1} achieves a strong baseline ($CC = 0.7425$), yet replacing real EEG with Gaussian noise produces a negligible $+0.23\%$ perturbation, and zeroing the visual image drops performance by $99.96\%$.

In the canonical replication of \textbf{Palazzo et al.} (\textit{IEEE TPAMI} 2021), zeroing visual features triggers a $+62.67\%$ error spike. In \textbf{Wang et al.} (\textit{CVPR} 2020), zeroing visual features induces a catastrophic $+228.35\%$ error explosion. In both cases, the shared decoder learns to exploit visual primitives, treating neural projections as uncalibrated offsets.

Conversely, the unimodal \textbf{EEGEyeNet} baseline demonstrates genuine neural sensitivity ($+114.10\%$ shift upon noise substitution). Because EEGEyeNet receives no visual input, the optimizer is strictly constrained to decode electrophysiological dynamics.

\subsection{Functional Cortical Lobe Ablation}
Table~\ref{tab:lobe_ablation} displays the impact of muting functional electrode clusters. Muting the primary visual area (Occipital lobe) results in a statistically negligible $0.06\%$ drop in $CC$. In biological vision, the P100 and N170 ERPs originate in occipito-temporal structures; an artificial network genuinely decoding cognitive visual attention must exhibit measurable sensitivity to these channels. The empirical invariance confirms that the neural pathway is bypassed.

\begin{table}[!t]
\caption{Cortical Lobe Channel Ablation (v1 Weights)}
\label{tab:lobe_ablation}
\centering
\footnotesize
\setlength{\tabcolsep}{3.5pt}
\begin{tabular}{llcc}
\toprule
\textbf{Cortical Region} & \textbf{Ablated Channels} & \textbf{Mean $CC$} & \textbf{$\Delta CC$ (\%)} \\
\midrule
Intact Montage & None (32 Channels) & 0.7425 & $0.00\%$ \\
Occipital Lobe & $O_1, O_2, O_z$ (Visual Area) & 0.7429 & $-0.06\%$ \\
Parietal Lobe & $P_3, P_4, P_z$ (Attention Control) & 0.7426 & $-0.01\%$ \\
Frontal Lobe & $F_3, F_4, F_z$ (Executive Control) & 0.7426 & $-0.02\%$ \\
Temporal/Central & $T_7, T_8, C_3, C_4$ (Sensory) & 0.7428 & $-0.04\%$ \\
\bottomrule
\end{tabular}
\end{table}

\begin{figure}[!t]
\centering
\includegraphics[width=0.48\textwidth]{outputs/figures/fig_spatial_disconnect.pdf}
\caption{The Spatial-Temporal Disconnect and Gradient Starvation Dynamics. (a) Spatial blur radius disparity between dipole volume-conducted scalp EEG ($>30\text{ mm}$) and convolutional visual feature maps ($<1\text{ mm}$). (b) Convergence disparity where visual gradients rapidly minimize loss, driving electrophysiological gradients to zero.}
\label{fig:spatial_disconnect}
\end{figure}

\section{Discussion: Why Algorithmic Mitigations Fail}

\subsection{Gate Collapse in Gated Residual Networks}
In unconstrained gated architectures (Version 3), the gate parameter $\text{Gate}(\mathbf{z}_E)$ rapidly converged to $\approx 10^{-3}$ for intermediate and shallow layers. Because visual features from ResNet-18 provide immediate loss minimization, the optimizer minimizes entropy by closing the residual gates against the high-entropy EEG stream.

\subsection{Degenerate Subspaces in Clamped Gating}
In Version 4, establishing a mathematical floor ($g_{\min} = 0.25$) prohibited the gate from closing. However, rather than learning meaningful EEG representations, the decoder weights $\theta_D$ adapted to project the forced neural subspace into a constant scalar bias. This collapsed predictive capacity ($CC = 0.1200$) while maintaining complete noise insensitivity ($\Delta CC = 0.00\%$), proving that forced gating without gradient re-balancing merely shifts the shortcut into the linear null-space of the decoder.

\subsection{The Spatial-Temporal Disconnect and Gradient Asymmetry}
As illustrated in Fig.~\ref{fig:spatial_disconnect}, scalp EEG provides millisecond temporal resolution but severe spatial blur ($>3\text{ cm}$) due to skull volume conduction. In free-viewing tasks on static imagery, fixation locations are dictated by localized 2D visual semantics (faces, high-contrast objects). Expecting temporal voltage fluctuations to refine spatial 2D coordinates beyond what a convolutional backbone provides is an ill-posed task under unconstrained loss landscapes. When both modalities are presented to a joint loss, the visual pathway's high signal-to-noise ratio in the spatial domain guarantees priority in parameter updates.

\subsection{The Limitation of Pareto Gradient Alignment (MMPareto)}
While MultiModal Pareto (MMPareto) projection successfully removes conflicting gradient components ($\langle g_{\text{main}}, g_{\text{aux}} \rangle < 0$), it assumes that both modalities reside on comparable manifold curvatures. When the visual gradient magnitude $\|g_V\|$ is several orders of magnitude larger than the neural gradient $\|g_E\|$, the projected update vector remains overwhelmingly collinear with the visual objective:
\begin{equation}
\lim_{\|g_V\| \gg \|g_E\|} \frac{g_{\text{Pareto}}}{\|g_{\text{Pareto}}\|} \approx \frac{g_V}{\|g_V\|}
\end{equation}
Consequently, orthogonalizing auxiliary objectives fails to resolve modality imbalance unless combined with explicit gradient rescaling or separate unimodal decoders.

\section{The Neuro-Visual Diagnostic Standard (NVDS)}
To prevent false multimodal claims in future literature, we formalize the \textbf{Neuro-Visual Diagnostic Standard (NVDS)}. A neuro-visual framework shall not be validated as multimodal unless it satisfies three mandatory criteria:
\begin{enumerate}
    \item \textbf{Waveform Sensitivity Criterion}: Replacing raw neural signals with Gaussian noise must induce a performance degradation exceeding a minimum threshold: $|\Delta \text{Metric}_{\text{noise}}| \ge 5.0\%$.
    \item \textbf{Trial Alignment Criterion}: Shuffling subject identities or temporal trial pairings must degrade performance: $|\Delta \text{Metric}_{\text{shuffle}}| \ge 2.0\%$.
    \item \textbf{Unimodal Retention Criterion}: Evaluating the model in the absence of visual stimuli must yield non-trivial predictive performance: $\text{Score}_{\text{zero-visual}} > \text{Chance Baseline}$.
\end{enumerate}

\section{Conclusion}
This study provides conclusive empirical evidence that reported high correlation metrics in multimodal neuro-visual saliency networks can be entirely illusory. Through systematic auditing of our own models and seminal published architectures (\textit{IEEE TPAMI}, \textit{CVPR}), we demonstrated that deep multimodal networks consistently succumb to the Greedy Learner Shortcut, bypassing neural signals in favor of pre-trained visual priors. Advanced algorithmic mitigations failed to halt this collapse. By establishing the Neuro-Visual Diagnostic Standard, we provide the neurocomputing community with a necessary verification protocol to distinguish genuine brain decoding from trivial unimodal shortcuts.

\bibliographystyle{IEEEtran}
\begin{thebibliography}{10}
\providecommand{\url}[1]{#1}
\csname url@samestyle\endcsname
\providecommand{\newblock}{\relax}
\providecommand{\bibinfo}[2]{#2}
\providecommand{\BIBentrySTDinterwordspacing}{\spaceskip=0pt\relax}
\providecommand{\BIBentryALTinterwordstretchfactor}{4}
\providecommand{\BIBentryALTinterwordspacing}{\spaceskip=\fontdimen2\font plus
\BIBentryALTinterwordstretchfactor\fontdimen3\font minus
  \fontdimen4\font\relax}
\providecommand{\BIBforeignlanguage}[2]{{%
\expandafter\ifx\csname l@#1\endcsname\relax
\typeout{** WARNING: IEEEtran.bst: No hyphenation pattern has been}%
\typeout{** loaded for the language `#1'. Using the pattern for}%
\typeout{** the default language instead.}%
\else
\language=\csname l@#1\endcsname
\fi
#2}}
\providecommand{\BIBdecl}{\relax}
\BIBdecl

\bibitem{riesenhuber1999hierarchical}
M.~Riesenhuber and T.~Poggio, ``Hierarchical models of object recognition in cortex,'' \emph{Nature Neuroscience}, vol.~2, no.~11, pp. 1019--1025, 1999.

\bibitem{itti1998model}
L.~Itti, C.~Koch, and E.~Niebur, ``A model of saliency-based visual attention for rapid scene analysis,'' \emph{IEEE Trans. Pattern Anal. Mach. Intell.}, vol.~20, no.~11, pp. 1254--1259, 1998.

\bibitem{cornia2018predicting}
M.~Cornia, L.~Baraldi, G.~Serra, and R.~Cucchiara, ``Predicting human eye fixations via an {LSTM}-based saliency attentive model,'' \emph{IEEE Trans. Image Process.}, vol.~27, no.~10, pp. 5142--5154, 2018.

\bibitem{kaushik2021eeg}
S.~Kaushik \emph{et~al.}, ``{EEG}-guided saliency: Integrating neural dynamics with visual feature extraction,'' \emph{NeuroImage}, vol. 240, p. 118356, 2021.

\bibitem{lanillos2020review}
P.~Lanillos \emph{et~al.}, ``A review on neural attention mechanisms in computer vision and cognitive robotics,'' \emph{IEEE Trans. Cogn. Dev. Syst.}, vol.~12, no.~4, pp. 642--658, 2020.

\bibitem{sugano2015self}
Y.~Sugano and A.~Bulling, ``Self-calibrating head-mounted eye trackers using {EEG} gaze markers,'' in \emph{Proc. ACM UIST}, 2015, pp. 439--448.

\bibitem{min2021decoding}
Z.~Min \emph{et~al.}, ``Decoding visual saliency from human brain signals during natural scene viewing,'' \emph{IEEE Trans. Neural Syst. Rehabil. Eng.}, vol.~29, pp. 2530--2539, 2021.

\bibitem{palazzo2020visual}
S.~Palazzo, F.~Rundo, S.~Battiato, D.~Giordano, and C.~Spampinato, ``Visual saliency detection guided by neural signals,'' in \emph{Proc. IEEE Int. Conf. Autom. Face Gesture Recognit. (FG)}, 2020, pp. 434--440.

\bibitem{palazzo2020decoding}
S.~Palazzo \emph{et~al.}, ``Decoding brain representations by multimodal learning of neural activity and visual features,'' \emph{IEEE Trans. Pattern Anal. Mach. Intell.}, vol.~43, no.~11, pp. 3833--3849, 2021.

\bibitem{wang2020makes}
W.~Wang, D.~Tran, and M.~Feiszli, ``What makes training multi-modal networks hard?'' in \emph{Proc. IEEE/CVF Conf. Comput. Vis. Pattern Recognit. (CVPR)}, 2020, pp. 12695--12704.

\bibitem{gifford2022alljoined}
K.~Gifford \emph{et~al.}, ``The {AllJoined} neuro-visual dataset: Paired high-density {EEG} and gaze tracking across {MS COCO} natural stimuli,'' \emph{Scientific Data}, vol.~9, no.~1, p. 412, 2022.

\bibitem{kastrati2021eegeyenet}
A.~Kastrati, M.~B. P{\l}omecka \emph{et~al.}, ``{EEGEyeNet}: A large-scale dataset and benchmark for simultaneous {EEG} and eye-tracking,'' in \emph{Adv. Neural Inf. Process. Syst. (NeurIPS) Track on Datasets and Benchmarks}, 2021.

\bibitem{perez2018film}
E.~Perez, F.~Strub, H.~de~Vries, V.~Dumoulin, and A.~Courville, ``{FiLM}: Visual reasoning with a general conditioning layer,'' in \emph{Proc. AAAI Conf. Artif. Intell.}, 2018, pp. 3942--3950.

\bibitem{vaswani2017attention}
A.~Vaswani \emph{et~al.}, ``Attention is all you need,'' in \emph{Adv. Neural Inf. Process. Syst. (NeurIPS)}, 2017, pp. 5998--6008.

\bibitem{peng2022balanced}
H.~Peng \emph{et~al.}, ``Balanced multimodal learning via on-the-fly gradient modulation,'' in \emph{Proc. IEEE/CVF Conf. Comput. Vis. Pattern Recognit. (CVPR)}, 2022, pp. 8238--8247.

\end{thebibliography}

\end{document}
